\chapter{Recerca prèvia}
\label{recerca previa}
\section{Marc normatiu batxillerat}
El batxillerat és una etapa postobligatòria que dura dos anys i, normalment, es fa després de l'ESO. Serveix perquè l'alumnat es prepari per entrar a la universitat o per fer cicles formatius de grau superior. A Catalunya hi ha quatre modalitats diferents: Ciències i Tecnologia, Humanitats i Ciències Socials, Arts i General. Tots tenen assignatures que són iguals per a tothom i d'altres que són les de cada itinerari.


Els que fan la modalitat de Ciències i Tecnologia acostumen a estudiar Matemàtiques II (MAT), i els de Ciències Socials acostumen a fer Matemàtiques Aplicades a les Ciències Socials II (MACS).

En el DECRET 171/2022, de 20 de setembre~\cite{curriculum-batxillerat-2022}, d'ordenació dels ensenyaments de batxillerat a l'article 13 punt 3 diu que un alumne no pot cursar Matemàtiques i Matemàtiques Aplicades a les Ciències Socials alhora, però, en canvi, sí que deixa que ho facin a les proves PAU, la qual cosa no és massa coherent. Aquest article s'ha modificat~\cite{curriculum-batxillerat-2026} i, a partir del curs 2026-2027, l'alumnat de qualsevol modalitat pot cursar les dues matèries alhora. Així, si l'horari del centre ho permet, es poden cursar totes dues presencialment, o bé cursar-ne una en línia a través de l'Institut Obert de Catalunya (IOC), i preparar-se així per als dos exàmens de les PAU.

\vspace*{1truecm}

\section{Què són les PAU?}
Les PAU són les Proves d'Accés a la Universitat. Hi ha una fase d'accés obligatòria i una altra voluntària per pujar nota. Cada estudiant decideix en quines matèries s'examina segons els estudis que vol fer després. Per això, algú que ha fet el batxillerat científic pot presentar-se a l'examen de MACS si això li pondera per a la carrera que vol cursar. Hi ha matèries que ponderen 0,2 i 0,1 en la nota final d'admissió, que es calcula:

$$\mbox{Nota d'admissió} = 0,6 \cdot \mbox{NMB} + 0,4 \cdot \mbox{NFA} + a \cdot \mbox{M1} + b \cdot \mbox{M2} $$.

On:
\begin{enumerate}
 \item NMB: nota mitjana de batxillerat.
 \item NFA: nota de la fase d'accés de les PAU (la part obligatòria).
 \item a i b: els coeficients de ponderació, que són 0,1 o 0,2 segons la matèria i el grau.
 \item M1 i M2: les dues millors notes de la fase específica (la part voluntària) entre les matèries que ponderen per al grau.
\end{enumerate}
El màxim és de $14$ punts: $10$ de la part general més fins a $4$ de la part específica ($2$ matèries $\cdot$ $10$ punts $\cdot$ $0,2$). Així, que realment val molt la pena fer l'examen de les matèries que ponderen $0,2$, ja que es nota molt la diferència de punts i més en graus on la nota de tall és alta, com en el cas de matemàtiques.

\vspace*{1truecm}

\section{Diferències entre els temaris de MAT i MACS}
La Generalitat publica els temaris de cada examen~\cite{Temari-mat,Temari-macs}. Els dos exàmens tenen gairebé la mateixa estructura, però el temari de cadascun és diferent, cada assignatura dona més importància a algunes coses. El temari del científic és molt més extens en la part d'algebra i anàlisi, per tant, un estudiant de MAT no ha de patir en aquesta part. En canvi, els estudiants de MAT no han fet estadística. A MACS, a més de la probabilitat que ja estudien al científic, també entra l'aproximació de la binomial per la normal, i a més a més, es fa la introducció al mostreig estadístic, l'estimació puntual de la mitjana, la proporció i la desviació típica, i intervals de confiança basats en la distribució normal.

% \clearpage

\section{Materials existents}
Vaig buscar si existia algun material (guia, llibre, vídeo...) pensat per a l'alumnat que ha cursat les MAT i a les PAU vol fer l'examen de MACS. Vaig mirar a biblioteques, en llibres de text, a YouTube i en plataformes d'apunts, però no vaig trobar cap recurs que expliqués només les diferències entre les dues assignatures. Hi ha llibres i vídeos que cobreixen tot el temari, però cap que compari directament els dos tipus de matemàtiques. A més, aquests materials sovint repeteixen conceptes que l'alumne de MAT ja domina i, en canvi, no aprofundeixen prou en els temes que li són nous.



%\section{Motivació del treball}
%Tota aquesta recerca neix de la voluntat d'estudiar el grau de Matemàtiques, on la nota de tall és molt alta i cal obtenir una nota molt alta a les matèries que ponderen: Matemàtiques, Matemàtiques aplicades a les ciències socials i Física. No hi ha materials específics per als quals, estudiants com jo, que han fet les Matemàtiques II, volen presentar-se a ambdos proves. Per això m'he proposat crear una guia clara i pràctica per ajudar a preparar aquest examen a companys que vinguin del batxillerat científic.
%He posat aquest info a introducció
